\subsection{Zeitverarbeitung und Synchronisation}
Die Zeitverarbeitung basiert auf der Nutzung eines \gls{ntp}-Servers zur Initialisierung sowie einer periodischen Nachsynchronisation zur Sicherstellung der Genauigkeit. 
Die aktuelle Zeit wird regelmäßig aus dem System gelesen und in globale Variablen überführt. 
Zusätzlich erfolgt täglich eine erneute Synchronisation, um Drift zu vermeiden. 
Der Ablauf ist in Abb.~\ref{fig:zeit} dargestellt.

\begin{figure}[H]
\centering
\begin{tikzpicture}[
    node distance=1.2cm,
    startstop/.style={rectangle, rounded corners, minimum width=4cm, draw=black},
    process/.style={rectangle, minimum width=4.5cm, draw=black},
    decision/.style={diamond, aspect=2.3, draw=black},
    io/.style={trapezium, trapezium left angle=70, trapezium right angle=110, minimum width=4.6cm, minimum height=0.8cm, text centered, draw=black},
    arrow/.style={->, thick}
]

\node (start) [startstop] {Zeit aktualisieren};
\node (read) [process, below=of start] {Systemzeit lesen};
\node (valid) [decision, below=of read] {Zeit gültig?};
\node (aktualisieren) [process, below=of valid] {Zeit aktualisieren};
\node (daily) [decision, below=of aktualisieren] {Mitternacht?};
\node (update) [process, below=of daily] {Tägliche Synchronisation mit NTP-Server};
\node (end) [startstop, below=of update] {Ende};

\draw [arrow] (start) -- (read);
\draw [arrow] (read) -- (valid);
\draw [arrow] (valid) -- node[right]{Ja} (aktualisieren);
\draw [arrow] (aktualisieren) -- (daily);
\draw [arrow] (daily) -- node[right]{Ja} (update);
\draw [arrow] (update) -- (end);
\draw [arrow] (valid.west) -- ++(-4,0) node[above]{Nein} |- (end.west);
\draw [arrow] (daily.west) -- ++(-3,0) node[above]{Nein} |- (end.west);

\end{tikzpicture}
\caption[Schematischer Programmablaufplan Zeitaktualisierung]{Schematischer Programmablaufplan zur Zeitaktualisierung und täglichen Synchronisation um Mitternacht.}
\label{fig:zeit}
\end{figure}